\section{Introduction}

Gaia parallaxes serve as foundational inputs to a broad range of astrophysical and cosmological analyses. Residual calibration and observation-geometry effects imprint coherent, sky-structured patterns in inferred parallaxes. Undiagnosed patterns risk propagation into downstream inference as if physical sky signals exist.

Quasars (QSOs) serve as an effectively zero-parallax reference population. The resulting parallax field constitutes a reference-frame residual surface whose spatial topology may reflect calibration error, contamination, or other unmodeled effects. The objective is characterization of that topology with emphasis on operational consequences for downstream use. Three conceptually separate axes are distinguished: (i) field amplitude (magnitude of sky-structured variation), (ii) concentration (fraction of weighted contribution carried by a minority of sky regions), and (iii) leverage (sensitivity of amplitude to exclusion of a small number of extreme regions).

A minimal, locked set of statistics and quality cuts is defined prior to inspection of results. The approach centers the contribution on diagnostic tooling rather than post-hoc rationalization.

The workflow quantifies whether a residual field appears (a) outlier-dominated or (b) broadly distributed via concentration and leverage. Validation of QSOs as an astrophysical zero-parallax frame is not performed. Mechanism identification (scan law, chromaticity, contamination, etc.) is not performed. Classification of calibration versus astrophysics is not performed.

